Thirty-nine designs were simulated in Stage~2 (Figs.~\ref{fig:ss2}--\ref{fig:ladder}; all at matched porosity $\phi=0.20\pm0.02$ and identical footprint unless stated, uniaxial stress along $x$). The observations that drove the hypotheses of Sec.~\ref{sec:hyp} are the following.

\paragraph{Orientation and load-path organisation dominate.} Slit pores parallel to the load give $\sigma_\mathrm{max}=25.4$\,N/m and $W=3.3$\,J/m$^2$, slits at $45^\circ$ 18.5\,N/m, slits perpendicular to the load 4.0\,N/m -- a factor 6.4 at identical porosity. Elliptical pores elongated along the load reach 19.0\,N/m against 8.6\,N/m across the load. In the strut lattices ($\phi=0.45$) the $0/90^\circ$ lattice with a continuous strut family along the load (11.5\,N/m) beats the triangulated $0/60/120^\circ$ lattice (7.4\,N/m, which relaxes strongly and fails in one avalanche) and the $30/90/150^\circ$ lattice without any load-parallel strut (6.0\,N/m); node connectivity (six- vs four-connected nodes) is thus much less important than whether straight through-going ligaments exist along the load (principles P8/P9).

\paragraph{Hierarchy does not add strength.} On the hierarchy ladder the coarse single-scale meshes are strongest (square pores on the 40\,\AA{} grid 17.0\,N/m, round pores with a 32\,\AA{} period 16.4\,N/m), the fine single-scale meshes weakest (11.7--13.3\,N/m for 16 and 12\,\AA{} periods) and all nested two-level meshes lie in between (12.7--14.3\,N/m): strength grows with vein width (4/8/12\,\AA: 13.2/12.8/13.7) and decreases with domain size (24/40/60\,\AA: 14.3/12.8/12.7); the three-level mesh (13.0\,N/m) equals the fine mesh. Failure strain and work to failure of the nested meshes (0.14--0.26, 1.0--1.5\,J/m$^2$) do not exceed those of the single-scale meshes (0.13--0.26, 0.9--1.9\,J/m$^2$). What hierarchy does change is the mode: the coarse square mesh accumulates harmless corner damage from 13.5\% strain and then snaps all load-parallel ligaments together at 17\%, whereas the nested meshes lose their fine level first (first damage at 11--13\%), retain 20--60\% of the peak load on the load-parallel veins and fail through crack bands that channel along columns adjacent to the load-\emph{perpendicular} veins (Sec.~\ref{sec:top}).

\paragraph{Disorder is a two-channel variable.} Pore-position jitter of 1/2/3\,\AA{} changes strength only mildly (12.7/11.8/11.1 versus 11.7\,N/m ordered) while lengthening the failure strain from 0.13 to 0.21--0.24 and multiplying the number of discrete damage events (6 $\to$ 17--22); pore-size dispersion of 15/30\% gives 12.0/11.3\,N/m with the same tail elongation. In the polygonal Voronoi networks the dependence is non-monotonic: random seeds 8.2--9.7\,N/m, 60\% regularised 14.4\,N/m (the best network, $W=1.6$\,J/m$^2$), fully regular 13.3\,N/m -- the random networks contain a few thin, misoriented ligaments that act as weak links and localise the damage ($L=0.40$--$0.42$).

\paragraph{Porosity and lattice orientation.} For the reference mesh the strength falls as 19.6/11.7/9.4/7.3\,N/m for $\phi=0.1/0.2/0.3/0.4$; the net-section strength $\sigma_\mathrm{max}/(1-\phi)$ itself decreases (21.7 $\to$ 12.2\,N/m), so thinner ligaments are intrinsically weaker in this model. Rotating the graphene lattice so that the armchair direction carries the load changes the mesh strength by only 3\% (11.3 vs 11.7\,N/m), whereas it changes the pristine sheet by 21\%: the architecture, not the lattice chirality, controls the porous sheets.

\paragraph{Gradients and flaw halos.} A pore-size gradient along or radial to the load (10.1 and 10.0\,N/m) is worse than the uniform mesh at equal porosity because the largest pores form the weakest section. Around a 20\,\AA{} hole, two rings of 5\,\AA{} pores (16.2\,N/m) cost less strength than the same porosity spread uniformly in the far field (13.2\,N/m), but both are below the bare hole (19.1\,N/m): the halo does not shield the flaw, it only localises the penalty of the added porosity.

\paragraph{Defect organisation changes the mode, not the strength.} Random and clustered vacancies give 27.5/27.2\,N/m at 1\%, 23.6/24.4 at 2\% and 22.4/23.4 at 3\%; in every pair the distributed defects fail progressively (6--9 damage steps, load retained) and the clustered ones in a single avalanche.

\paragraph{A single descriptor explains most of the strength variation.} Across all 39 designs the strength is best organised by the minimum load-bearing solid fraction across the loading direction (Fig.~\ref{fig:maps}); the ratio $\sigma_\mathrm{max}/\mathrm{minSF}$ (a net-section strength) is 29--32\,N/m for straight, load-parallel ligaments (slits, square and round coarse pores, load-elongated ellipses), 17--19\,N/m for fine round-pore meshes with ligaments of 10--13\,\AA{}, 22\,N/m for the 32\,\AA{} round-pore mesh, and 12--15\,N/m for random networks. Neither porosity, ligament width, edge-atom fraction, anisotropy index nor tortuosity alone orders the data as well.
